\section{Introduction}

Multi-agent systems built on autonomous spawning architectures face a fundamental fragility that the literature rarely names directly: \textbf{autonomous drift}. An agent spawns a child to handle a subtask; the child spawns another; permissions attenuate, context narrows, and somewhere several generations down the lineage, an agent begins operating on incentives, goals, or factual grounding that diverge — sometimes subtly, sometimes catastrophically — from those of its ancestors. By the time this divergence surfaces in observable behavior, the causal chain has dissolved into a tangle of implicit delegation decisions no individual agent can reconstruct. This is not a failure of alignment or instruction quality. It is a structural vulnerability endemic to spawning architectures that lack a persistent, integrity-preserving lineage mechanism.

\subsection{The Drift Problem}

The problem manifests when spawn chains exceed the information horizon of any single monitoring agent. Consider a scenario in the BX3 substrate: a \texttt{cascade} trigger initiates a spawning chain in response to a market anomaly. The root agent (\texttt{Gateway}) delegates to Agent Alpha, which delegates to Agent Beta, which delegates to Agent Gamma. Each delegation inherits the parent's authorization context and goal frame, but none of these inherit a verifiable commitment mechanism — the child is free to reinterpret the parent's intent as its own context evolves. If Gamma's reasoning process introduces a factual error or a goal reframing, neither Beta nor Alpha has the structural means to detect it before the error propagates into downstream actions such as cascade firing or settlement execution.

This is distinct from the classic Byzantine fault problem. Byzantine faults assume adversarial or arbitrary node behavior; autonomous drift requires no malice. It arises from the perfectly legitimate capacity of each agent to adapt its reasoning to local context, which, absent a lineage integrity mechanism, degrades the chain's collective coherence in ways that are neither detectable nor reversible.

\subsection{Why Existing Approaches Fail}

The two dominant mitigations in current multi-agent frameworks — TTL-based spawn limits and depth-limited trees — address the symptoms of drift without touching its cause.

\textbf{TTL-based termination} caps the lifetime of any spawned agent, forcing periodic re-evaluation of whether a spawned process should continue. While TTL prevents unbounded chains, it does nothing to preserve the semantic integrity of what the chain has done or intended. A short-lived agent that has already emitted a cascade trigger or modified shared state has left a footprint that TTL does not erase. Moreover, TTL is a blunt instrument: a long TTL permits deep, drifting chains; a short TTL interrupts legitimate long-running tasks arbitrarily.

\textbf{Depth-limited trees} restrict how many generations of spawning are permitted, typically by assigning a depth budget that decrements with each spawn. This limits exposure but introduces a new problem: the depth limit is a global parameter, and its correct value depends on the complexity of the task tree, which the designer cannot know \textit{a priori}. Setting it too low fragments legitimate task decomposition; setting it too high permits drift to accumulate across too many generations. Neither approach provides the system with a way to verify that a child agent's behavior is consistent with its ancestor's intent — they merely bound the blast radius.

A third class of approaches, checkpoint-and-revert mechanisms, attempt to make lineage auditable by snapshotting agent state at spawn time. These are operationally expensive and introduce consistency hazards when a parent continues operating while a child's snapshot is being evaluated. None of these approaches address the core issue: the absence of an immutable, publicly verifiable lineage record that every agent in the chain can reference.

\subsection{Immutable Parent Pointer Chains: The Structural Fix}

The core insight of this paper is that autonomous drift is not a behavioral problem — it is a structural one. It can be made \textbf{structurally impossible} by replacing the implicit, mutable delegation context of spawn chains with an \textbf{immutable parent pointer chain}. In the BX3 formulation, every spawned agent receives at instantiation a cryptographic commitment to its parent's canonical identifier and intent fingerprint. This commitment is stored in the agent's immutable header and is propagated — not copied — to any agent it in turn spawns.

The key property is immutability: a parent pointer in this architecture \textit{cannot be rewritten}. An agent may terminate, suspend, or be superseded, but its identity and its lineage commitment persist in the substrate. Any agent that queries its own lineage receives a tamper-evident chain from the current agent back to the root spawn — the \texttt{Gateway} agent in BX3 terminology. If any link in the chain is missing or contains an intent mismatch, this is detected structurally, not behaviorally.

This transforms drift from a runtime coherence problem into a compile-time impossibility. If every agent's lineage is structurally bounded by an immutable pointer chain, then an agent cannot drift in the relevant sense: it may fail, malfunction, or receive corrupted context, but it cannot sustain a goal or factual divergence that is invisible to the lineage verification mechanism, because the lineage verification mechanism is built into the chain itself.

\subsection{Formal Definition of Drift}

We define \textbf{autonomous drift} formally as follows.

Let $A_0$ be a root agent with intent vector $\mathbf{I}(A_0)$ and factual grounding $\mathbf{F}(A_0)$. At spawn, agent $A_i$ inherits from parent $A_{i-1}$ an intent inheritance bundle $(\mathbf{I}(A_{i-1}), \mathbf{F}(A_{i-1}))$. Agent $A_i$ is said to \textbf{drift} if, after $k \geq 0$ reasoning cycles within its own context, its maintained intent vector $\mathbf{I}'(A_i)$ and factual grounding $\mathbf{F}'(A_i)$ satisfy:

\[
\exists\, j \in \{0, \dots, i-1\} \text{ such that } \text{distance}(\mathbf{I}'(A_i), \mathbf{I}(A_j)) > \tau_I \quad \text{or} \quad \text{distance}(\mathbf{F}'(A_i), \mathbf{F}(A_j)) > \tau_F,
\]

where $\tau_I$ and $\tau_F$ are threshold constants for intent and factual divergence respectively, and the divergence is \textbf{undetectable by any agent in the chain} at the time of its occurrence.

In a system with an immutable parent pointer chain, the parent pointer of $A_i$ encodes a commitment to $\mathbf{I}(A_{i-1})$ and $\mathbf{F}(A_{i-1})$. Any divergence beyond $\tau$ produces a lineage verification failure that propagates a fault signal to any querying agent, satisfying the detectability condition. This eliminates the ``undetectable'' qualifier, which is the critical property that makes drift dangerous.

\subsection{Paper Roadmap}

The remainder of this paper is organized as follows. Section~\ref{section:architecture} presents the immutable parent pointer architecture in detail, including the cryptographic commitment scheme, the substrate integration points for the BX3 substrate, and the formal properties of the pointer chain. Section~\ref{section:verification} introduces the lineage verification protocol and the drift detection criteria that the substrate enforces structurally. Section~\ref{section:experiments} describes experimental evaluations of the architecture against baseline spawning systems on tasks of varying depth and complexity. Section~\ref{section:discussion} discusses limitations, edge cases involving concurrent spawning and memory-partitioned agents, and the relationship between immutable lineage and existing alignment frameworks. Section~\ref{section:conclusion} concludes with implications for the design of production multi-agent systems and a set of open problems.

% Placeholder for subsequent sections
\section{Architecture}
\label{section:architecture}
% Content to follow.

\section{Lineage Verification and Drift Detection}
\label{section:verification}
% Content to follow.

\section{Experimental Evaluation}
\label{section:experiments}
% Content to follow.

\section{Discussion}
\label{section:discussion}
% Content to follow.

\section{Conclusion}
\label{section:conclusion}
% Content to follow.